% ---------------------------------------------------------------------
\section{The Technology Stack}
\label{sec:techstack}
% ---------------------------------------------------------------------

Everything here \emph{extends} the original Go~+~Rust~+~llama.cpp design; nothing
is rewritten. The one mindset shift: \textbf{the controller owns reliability,
not llama.cpp}.

\subsection{Component by component}

\field{Controller --- Go \Keep} The coordinator ``brain''. Handles the HTTP API
(user sends prompts), the task queue, worker heartbeats (detect a dead laptop),
result collection, and --- the novel part --- the \textbf{failover logic}:
``node died → who is alive and has capacity → reassign its layers → re-run the
request''. Go is well suited to this and the heartbeat/queue pieces already exist.

\field{Worker agent --- Rust \Keep} Runs on each laptop. Registers with the
controller, sends heartbeats (``alive; here is my CPU/RAM/battery''), loads its
assigned slice of the model, runs its layers via llama.cpp, and passes
activations to the next worker. Rust compiles to Windows, Linux, and macOS.

\field{Inference engine --- llama.cpp \Keep (usage changed)}
\begin{itemize}[leftmargin=1.4em,nosep]
  \item \Drop\ \textbf{Do NOT use llama.cpp's built-in RPC mode.} It crashes on
  node loss (\VER, \S\ref{sec:landscape}); building on it inherits the exact
  failure we are beating.
  \item \Keep\ \textbf{Use llama.cpp as a local engine on each worker}, with the
  Go controller in charge of coordination and recovery. We route \emph{around}
  llama.cpp's weak part. This single decision is what makes the project possible.
\end{itemize}

\field{Model --- Meta Llama, small, GGUF \Keep} A small open-weights Llama (e.g.
Llama~3.2 1B or 3B) in GGUF, quantized (e.g. Q4). Small = fast load, fits weak
laptops, snappy demo. Satisfies the hackathon's ``use a Meta model'' requirement.
\Drop\ 70B models --- slow demo, and prima.cpp wins that contest anyway.

\field{Networking --- TCP/gRPC over Wi-Fi \Keep} Workers $\leftrightarrow$
controller over plain TCP/JSON to start (upgrade to gRPC only if needed). Wi-Fi
is both the demo environment and the story (``real laptops, real Wi-Fi'').

\field{Model split --- pipeline (by layers) \Keep} Split by layers (worker~1 =
layers 0--8, worker~2 = 9--16, \ldots). For recovery, a survivor picks up the
dead worker's layers.

\field{Recovery method --- request-level re-route \Keep (the core)}
When a worker dies mid-request, the controller reassigns its layers to a survivor
and \textbf{re-runs the current request} from the start of generation. The user
sees a short pause, then continuation. \Drop\ mid-token activation handoff --- too
hard for the time window.

\field{Code written with --- Kiro \Keep} The Go controller, the Rust worker, the
failover handling, and the glue are written with Kiro.

\subsection{Summary table}
\begin{table}[h]
\centering\footnotesize
\caption{HyperPod technology stack}
\label{tab:stack}
\begin{tabularx}{\linewidth}{L{0.34\linewidth}Y}
\toprule
Piece & Choice \\
\midrule
Controller / failover brain & Go (keep, extend) \\
Worker agent per laptop & Rust (keep) \\
Inference engine & llama.cpp --- local mode, NOT RPC (keep, change usage) \\
Model & Meta Llama 3.2 1B/3B, GGUF, quantized (add) \\
Networking & TCP/JSON then gRPC over Wi-Fi (add) \\
Model split & by layers (pipeline) (add) \\
Recovery method & request-level re-route, controller-driven (core) \\
Code written with & Kiro \\
\bottomrule
\end{tabularx}
\end{table}
